\chapter{A deposit-policy study from rollout to funding value}\label{ch:deposits}
\section{The owner's decision and the financial account}
Suppose the group funding committee is considering a temporary savings-rate
enhancement. The retail business expects retained deposits; treasury values
the funding; finance pays interest on qualifying balances; and the customer
receives a higher rate. The relevant bank decision is whether the incremental
funding value exceeds the incremental interest and operating costs over a
specified horizon. An effect on account balances is one input to that
decision, not its answer.

The committee needs three quantities kept separate: balances in the promoted
product, total external funding retained by the group, and incremental
economic contribution. An internal transfer from another group account can
increase the first with little change in the second. Paying a higher rate
on incumbent balances can make the third negative even when both balance
measures rise. This decomposition should be visible in the first report.

The empirical literature suggests mechanisms and designs. The deposits
channel in \citet{drechsler2017} connects monetary policy, local deposit
market power, funding flows and lending. Its within-bank branch comparison
uses geographic concentration while controlling bank-specific opportunities;
its interpretation relies on banks being able to move funding between
branches. The paper's subsequent lending comparisons operate at different
aggregation levels. It does not supply the effect of a new bank's account-
level retention offer. The product should make that distinction explicit
when an analyst imports a mechanism from a published study.

\section{The data contract before estimation}
The proposed bank implementation begins with a fixed eligible cohort before
announcement. It preserves branch and household identity as of that date,
the original assignment, the announced terms, delivery, take-up and later
closure. It then reconstructs daily balances, interest, external flows and
internal transfer links under a stated accounting cutoff. A corrected
transaction changes a later readout while leaving the original information
set and assignment intact.

\begin{table}[htbp]\centering\small
\begin{tabularx}{\textwidth}{P{31mm}XX}
\toprule Object & Required meaning & Consequence if absent\\\midrule
Eligibility & Population defined before announcement & Survivor-only analysis changes the target\\
Assignment & Offer terms, duration, branch/cohort and date & Acceptance cannot substitute for assignment\\
Balance & Stock, currency, product and economic date & Daily and monthly-average balances are not interchangeable\\
Flows & External source or linked internal transfer & Group funding cannot be inferred from one product's growth\\
Interest & Actual accrual and promotional increment & Repricing of incumbent funds is omitted\\
Funding value & Dated marginal funding convention & Principal may be mistaken for contribution\\
Observation & First available date and correction lineage & Retrospective leakage changes the study information set\\
\bottomrule
\end{tabularx}
\caption{The proposed deposit study's minimum economic objects. A bank may
support a narrower balance study before all funding objects are available.}
\end{table}

Households moving between branches require a declared attribution rule.
Freezing baseline assignment preserves an intention-to-treat interpretation;
following current branch membership can select on the policy response.
Cross-branch displacement also challenges untreated donors. A group-wide
study can use household-level outcomes while clustering by randomized or
rollout branch, provided the mapping is stable and spillovers are addressed.

\section{A complete synthetic rollout}
The executed study observes 96 branches for 30 months. Twenty-four branches
adopt in each of months 12, 18 and 24, with 24 never adopting. Months are
indexed from zero. Cohort labels are fixed and branch disturbances are
independent of them. Untreated balances, in millions of dollars, are
\begin{align}
 B_{it}(0)&=30+\alpha_i+0.07t+0.4\sin(t/3)+\epsilon_{it},\\
 \alpha_i&\sim N(0,3^2),\qquad
 \epsilon_{it}=0.65\epsilon_{i,t-1}+\eta_{it},\quad
 \eta_{it}\sim N(0,0.28^2),
 \label{eq:deposit-dgp}
\end{align}
with the first disturbance drawn at the innovation scale. For cohort $g$
and event month $k=t-g\geq0$, the balance effect is
\[
 \tau_{gk}=\{0.85+0.025(g-12)\}(1+0.14k).
\]
This law imposes untreated parallel trends while allowing heterogeneous,
growing effects. It is designed to compare targets and estimators rather
than simulate every feature of deposit demand.

For every supported $(g,t)$, estimate the change from $g-1$ among that
cohort minus the same change among never-treated branches. Average all
treated cohort-month cells equally. Because all cohorts have 24 branches,
this equals an equal-weight average over treated branch-months. It gives
early cohorts more total weight because they contribute more treated
months. A bank wanting an equal-cohort or balanced-event-window target
would use different declared weights.

The challenger fits a single treatment indicator with branch and month
effects. It estimates a different weighted combination when effects evolve
with adoption time, as explained in Chapter~\ref{ch:panels}. The study
also calculates a known reference by comparing the generated potential
outcomes for the same branch-months. Those potential outcomes are available
only to the synthetic evaluator.

\section{Converting balances into contribution}
The teaching account assumes an annual net funding margin of $2.5\%$
on balances before the promotional increment, a 25-basis-point promotional
rate increase on all treated balances, and USD\,1,800 monthly program
cost per treated branch. Under the stipulated monthly convention,
\begin{align}
 C_{it}(1)&=10^6 B_{it}(1)\frac{0.025}{12}
           -10^6 B_{it}(1)\frac{0.0025}{12}-1800,\\
 C_{it}(0)&=10^6 B_{it}(0)\frac{0.025}{12}.
 \label{eq:deposit-cash}
\end{align}
Subtracting gives
\begin{equation}
 \Delta C_{it}=10^6\frac{0.025-0.0025}{12}\Delta B_{it}
       -10^6\frac{0.0025}{12}B_{it}(0)-1800.
 \label{eq:deposit-hurdle}
\end{equation}
The second term is the cost of repricing balances that would have remained.
It is easy to omit when a business case values only additional deposits.

At an untreated balance of USD\,30 million, break-even requires
\[
 \Delta B>
 \frac{0.0025(30)+12(1800)/10^6}{0.025-0.0025}
 \simeq 4.293\ \hbox{million dollars}.
\]
The hurdle follows directly from the financial account. It changes with
rate duration, funding value, churn, servicing and the share of incumbent
balances receiving the offer. The numerical study holds those terms fixed
so that the effect of repricing is transparent.

\section{Fitted results and uncertainty}
The same cohort-time estimator is applied separately to balances and
observed contribution. Whole branches are resampled within each adoption
cohort in 499 bootstrap samples, retaining each trajectory and the fixed
cohort counts. The percentile intervals condition on this rollout design
and the supplied common macro path. They are not a repeated-sample
coverage qualification for a bank rollout.

\begin{table}[htbp]\centering\small
\begin{tabular}{@{}lrrrr@{}}
\toprule
Outcome & True effect & Cohort ATT & 95\% interval & TWFE\\
\midrule
Balances (USD m) & 1.780 & 1.689 & [1.564, 1.816] & 1.442\\
Contribution (USD) & -5030.635 & -5218.875 & [-5498.095, -4935.700] & -5607.561\\
\bottomrule
\end{tabular}

\caption{One fixed synthetic rollout. The reference is the realized
finite-population potential-outcome contrast; intervals resample complete
branches. Positive balance effects coexist with negative contribution.}
\label{tab:deposit-results}
\end{table}

\begin{figure}[htbp]\centering
\includegraphics[width=\textwidth]{derivations/deposit_study.pdf}
\caption{Cohort balance responses and the funding-contribution comparison.
The differing event profiles matter for aggregation; their positive signs
do not establish a profitable promotion.}
\label{fig:deposit-study}
\end{figure}

The contribution estimate is approximately USD\,$-5,219$ per treated
branch-month, with interval $[-5,498,-4,936]$. Under the stipulated design
and account, the rollout loses money despite retaining balances. The
appropriate decision is to reject this broad offer or investigate a
different, explicitly evaluated policy. Selecting a more favorable
balance estimator would not remove the cost of repricing incumbent funds.

The balance and contribution estimates share outcomes and branches.
For a portfolio of more detailed cash components, the implementation
retains their joint bootstrap or influence contributions. It should not
form a purported interval for total contribution by adding separately
calculated marginal endpoints. The simple account here allows contribution
to be observed directly, avoiding an unnecessary two-stage uncertainty
approximation.

\section{A targeted alternative is a new study}
A natural response is to offer the enhancement only to households likely
to withdraw external funds. That may reduce repricing cost, but it changes
the intervention and target population. A withdrawal-risk model predicts
which customers leave under historical practice. It does not necessarily
predict whose withdrawal is preventable by the offer. High-risk customers
may be insensitive to price; low-risk customers may accept a subsidy
without changing their behavior.

The appropriate targeting comparison uses conditional effects on the
group funding account and evaluates the whole targeting rule on held-out
or randomized data. It includes contact costs, take-up, rate duration and
customer constraints. The existing static example demonstrates why a
known positive average effect does not establish that one targeted policy
beats another at the declared simultaneous confidence level.

For repeat offers, the problem becomes longitudinal. A first offer changes
balances and may teach customers to wait for future offers. Selection for
the next offer then depends on a variable affected by the first.
The sequential methods in Chapter~\ref{ch:longitudinal} are needed to
compare sustained strategies. A one-period uplift score should not be
repeated indefinitely under an assumption of unchanged effects.

\section{What can transfer from the literature}
\citet{drechsler2017} supports a mechanism linking market power,
deposit pricing and transmission, with specific within-bank and local
comparisons. \citet{basten2026} investigates relationship stickiness and
cross-selling using mover and prior-location variation. Both suggest
heterogeneity worth measuring. Neither establishes that a bank's current
customers share the same elasticities, product alternatives or information.

Refinancing evidence adds a complementary lesson. In
\citet{agarwal2020}, eligibility for HARP provides a policy comparison;
the paper distinguishes a program effect from the selected experience
of borrowers who refinance and studies competitive frictions in pass-through.
This is analogous to keeping an assigned retention offer distinct from
acceptance. \citet{beraja2019} further explains why local responses do
not settle aggregate policy effects. A bank-wide price change can alter
competitor behavior and internal funding substitution beyond the local
partial-equilibrium study.

The platform should consequently store an evidence-transfer statement:
which mechanism or design comes from the paper, which coefficient is
estimated on owned data, and which equilibrium or stability assumptions
remain scenario inputs. Literature informs construction without becoming
an unexamined parameter database.

\section{The reusable application and acceptance test}
The deposit application exports a study specification, reconstructed cohort,
balance/flow/contribution definitions, cohort-time effects, covariance,
target weights, trend sensitivity and a funding report. A FinMarketing
consumer can use the supported causal response in its own financial
workflow if the population, intervention and account match. A treasury
consumer can use the funding contrast while retaining its transfer-price
and liquidity assumptions.

Bank acceptance would require reconciliation to authoritative balances
and interest, verified external-versus-internal flows, a defensible rollout
or experimental design, explicit spillover analysis, adequate independent
branches and outcome maturity. The equivalent-task comparator is a bank
economist building the same study with existing panel software and the
warehouse. The product earns its role if it reconstructs and updates that
study more reliably and cheaply while preserving the same scientific
answer, including this kind of unfavorable result.
